% GENERATED FILE. Edit canonical lessons, then run npm run generate:books.
% canonical-source-hash: fnv1a64:0d3a04a7f35a5976
% canonical-lessons: UR-C437-jhari, UR-C437-saya, UR-C437-ghana, UR-C437-kai, UR-C437-ugna

\chapter{Bushes and Shade}
\label{ch:ur-bushes-and-shade}
\hlchaptermodality{437}

\begin{chapteropening}
\textbf{By the end of this chapter:} \emph{I can describe thick growth, shade and green places.}

Complete the last lesson of chapter 437: I can describe thick growth, shade and green places.
\end{chapteropening}

\section[\texorpdfstring{\ur{جھاڑی} (jhāṛī)}{جھاڑی (jhāṛī)}]{\ur{جھاڑی} (jhāṛī) — a bush}

\label{lesson:UR-C437-jhari}



\begin{quote}
Before the new one: say the Urdu for a petal, then the Urdu for to wilt.
\end{quote}

\subsection*{You'll want to know: \ur{جھاڑی}}
\textbf{\ur{جھاڑی}} — \emph{jhāṛī} — \textquotedblleft{}a bush\textquotedblright{}.

\begin{quote}
\ur{سڑک} \ur{کے} \ur{کنارے} \ur{جھاڑیاں} \ur{ہیں۔}
\end{quote}

\begin{quote}
\emph{saṛak ke kināre jhāṛiyāṅ haiṅ.}
\end{quote}

\begin{quote}
There are bushes along the road.
\end{quote}

Feminine.

\begin{cousinweb}
\textbf{\ur{جھاڑی}} (\emph{jhāṛī}) is inherited: CDIAL lists the Hindi cousins \emph{jhāṛ} and \emph{jhāṛī}, \textquotedblleft{}a bush\textquotedblright{}, under \emph{jhāṭa}, \textquotedblleft{}a forest, an arbour\textquotedblright{} (CDIAL 5362).
\end{cousinweb}

\begin{grammarlens}[title={What you've built}]
One more word for getting around.
\end{grammarlens}

\subsection*{Your turn}
\begin{itemize}
\raggedright
  \item \emph{Say it:} \emph{jhāṛī}
  \item \emph{Say it:} \emph{jhāṛī}, once more
  \item \emph{Say it:} \emph{saṛak ke kināre jhāṛiyāṅ haiṅ.} — There are bushes along the road.
  \item \emph{Recall:} say the Urdu for a petal, then the Urdu for to wilt, then say \emph{jhāṛī} again
\end{itemize}

\begin{tcolorbox}[breakable,colback=teal!4,colframe=teal!35!black,title={Before you move on}]
What does \ur{جھاڑی} mean? (\textquotedblleft{}A bush\textquotedblright{}.) Say it once more.
\end{tcolorbox}
\section[\texorpdfstring{\ur{سایہ} (sāyā)}{سایہ (sāyā)}]{\ur{سایہ} (sāyā) — shade; a shadow}

\label{lesson:UR-C437-saya}



\begin{quote}
Before the new one: say the Urdu for to wilt, then the Urdu for a bush.
\end{quote}

\subsection*{You'll want to know: \ur{سایہ}}
\textbf{\ur{سایہ}} — \emph{sāyā} — \textquotedblleft{}shade; a shadow\textquotedblright{}.

\begin{quote}
\ur{گرمی} \ur{میں} \ur{ہم} \ur{درخت} \ur{کے} \ur{سائے} \ur{میں} \ur{بیٹھے۔}
\end{quote}

\begin{quote}
\emph{garmī meṅ ham darakht ke sāe meṅ baiṭhe.}
\end{quote}

\begin{quote}
In the heat we sat in the shade of a tree.
\end{quote}

Masculine.

\begin{cousinweb}
\textbf{\ur{سایہ}} (\emph{sāyā}) is usually traced to Persian \emph{sāya}, \textquotedblleft{}shade, a shadow\textquotedblright{}.
\end{cousinweb}

\begin{grammarlens}[title={What you've built}]
One more word for getting around.
\end{grammarlens}

\subsection*{Your turn}
\begin{itemize}
\raggedright
  \item \emph{Say it:} \emph{sāyā}
  \item \emph{Say it:} \emph{sāyā}, once more
  \item \emph{Say it:} \emph{garmī meṅ ham darakht ke sāe meṅ baiṭhe.} — In the heat we sat in the shade of a tree.
  \item \emph{Recall:} say the Urdu for to wilt, then the Urdu for a bush, then say \emph{sāyā} again
\end{itemize}

\begin{tcolorbox}[breakable,colback=teal!4,colframe=teal!35!black,title={Before you move on}]
What does \ur{سایہ} mean? (\textquotedblleft{}Shade; a shadow\textquotedblright{}.) Say it once more.
\end{tcolorbox}
\section[\texorpdfstring{\ur{گھنا} (ghanā)}{گھنا (ghanā)}]{\ur{گھنا} (ghanā) — dense; thick}

\label{lesson:UR-C437-ghana}



\begin{quote}
Before the new one: say the Urdu for a bush, then the Urdu for shade.
\end{quote}

\subsection*{You'll want to know: \ur{گھنا}}
\textbf{\ur{گھنا}} — \emph{ghanā} — \textquotedblleft{}dense; thick\textquotedblright{}.

\begin{quote}
\ur{یہ} \ur{جنگل} \ur{گھنا} \ur{ہے۔}
\end{quote}

\begin{quote}
\emph{yih jangal ghanā hai.}
\end{quote}

\begin{quote}
This forest is dense.
\end{quote}

\begin{cousinweb}
\textbf{\ur{گھنا}} (\emph{ghanā}) is inherited: CDIAL lists the Hindi cousin \emph{ghanā}, \textquotedblleft{}dense\textquotedblright{}, under Sanskrit \emph{ghaná}, \textquotedblleft{}compact, dense\textquotedblright{} (CDIAL 4424).
\end{cousinweb}

\begin{grammarlens}[title={What you've built}]
One more word for getting around.
\end{grammarlens}

\subsection*{Your turn}
\begin{itemize}
\raggedright
  \item \emph{Say it:} \emph{ghanā}
  \item \emph{Say it:} \emph{ghanā}, once more
  \item \emph{Say it:} \emph{yih jangal ghanā hai.} — This forest is dense.
  \item \emph{Recall:} say the Urdu for a bush, then the Urdu for shade, then say \emph{ghanā} again
\end{itemize}

\begin{tcolorbox}[breakable,colback=teal!4,colframe=teal!35!black,title={Before you move on}]
What does \ur{گھنا} mean? (\textquotedblleft{}Dense; thick\textquotedblright{}.) Say it once more.
\end{tcolorbox}
\section[\texorpdfstring{\ur{کائی} (kāī)}{کائی (kāī)}]{\ur{کائی} (kāī) — moss; green scum (on water)}

\label{lesson:UR-C437-kai}



\begin{quote}
Before the new one: say the Urdu for shade, then the Urdu for dense.
\end{quote}

\subsection*{You'll want to know: \ur{کائی}}
\textbf{\ur{کائی}} — \emph{kāī} — \textquotedblleft{}moss; green scum (on water)\textquotedblright{}.

\begin{quote}
\ur{تالاب} \ur{کے} \ur{پانی} \ur{پر} \ur{کائی} \ur{ہے۔}
\end{quote}

\begin{quote}
\emph{tālāb ke pānī par kāī hai.}
\end{quote}

\begin{quote}
There is green scum on the pond.
\end{quote}

Feminine.

\begin{cousinweb}
\textbf{\ur{کائی}} (\emph{kāī}) is inherited: CDIAL lists the Hindi cousin \emph{kāī}, \textquotedblleft{}green scum on still water\textquotedblright{}, under \emph{kāvikā}, a form it reconstructs (CDIAL 3109).
\end{cousinweb}

\begin{grammarlens}[title={What you've built}]
One more word for getting around.
\end{grammarlens}

\subsection*{Your turn}
\begin{itemize}
\raggedright
  \item \emph{Say it:} \emph{kāī}
  \item \emph{Say it:} \emph{kāī}, once more
  \item \emph{Say it:} \emph{tālāb ke pānī par kāī hai.} — There is green scum on the pond.
  \item \emph{Recall:} say the Urdu for shade, then the Urdu for dense, then say \emph{kāī} again
\end{itemize}

\begin{tcolorbox}[breakable,colback=teal!4,colframe=teal!35!black,title={Before you move on}]
What does \ur{کائی} mean? (\textquotedblleft{}Moss; green scum (on water)\textquotedblright{}.) Say it once more.
\end{tcolorbox}
\section[\texorpdfstring{\ur{اگنا} (ugnā)}{اگنا (ugnā)}]{\ur{اگنا} (ugnā) — to come up (of plants); to grow}

\label{lesson:UR-C437-ugna}



\begin{quote}
Before the new one: say the Urdu for dense, then the Urdu for moss.
\end{quote}

\subsection*{You'll want to know: \ur{اگنا}}
\textbf{\ur{اگنا}} — \emph{ugnā} — \textquotedblleft{}to come up (of plants); to grow\textquotedblright{}.

\begin{quote}
\ur{بارش} \ur{کے} \ur{بعد} \ur{کھیتوں} \ur{میں} \ur{گھاس} \ur{اگی۔}
\end{quote}

\begin{quote}
\emph{bārish ke bād khetoṅ meṅ ghās ugī.}
\end{quote}

\begin{quote}
After the rain grass came up in the fields.
\end{quote}

Past (it …): \ur{وہ} \ur{پودا} \ur{اگا}, \emph{voh paudā ugā}.

\begin{cousinweb}
\textbf{\ur{اگنا}} (\emph{ugnā}) is inherited: CDIAL lists the Hindi cousin \emph{ugānā}, \textquotedblleft{}to cause to spring up\textquotedblright{}, under \emph{udgāti}, \textquotedblleft{}rises, comes forth\textquotedblright{}, a form it reconstructs (CDIAL 1954); \textbf{\ur{اگانا}}, \textquotedblleft{}to grow (crops)\textquotedblright{}, is the partner for making things come up.
\end{cousinweb}

\begin{grammarlens}[title={What you've built}]
One more word for getting around.
\end{grammarlens}

\subsection*{Your turn}
\begin{itemize}
\raggedright
  \item \emph{Say it:} \emph{ugnā}
  \item \emph{Say it:} \emph{ugnā}, once more
  \item \emph{Say it:} \emph{bārish ke bād khetoṅ meṅ ghās ugī.} — After the rain grass came up in the fields.
  \item \emph{Recall:} say the Urdu for dense, then the Urdu for moss, then say \emph{ugnā} again
\end{itemize}

\begin{tcolorbox}[breakable,colback=teal!4,colframe=teal!35!black,title={Before you move on}]
What does \ur{اگنا} mean? (\textquotedblleft{}To come up (of plants); to grow\textquotedblright{}.) Say it once more.
\end{tcolorbox}
